\begin{textsample}{POS Dim 1 – sc – Score 48.00 – sc/sc\_em\_31.txt}  \label{ex:f1_pos_011}
3.2 COMPETÊNCIAS, HABILIDADES E CONCEITOS ESTRUTURANTES NA ÁREA DE CIÊNCIAS HUMANAS E SOCIAIS APLICADAS Para \textbf{que} haja aprendizagens significativas dos(das) estudantes do ensino médio, a proposta da Base Nacional Comum Curricular ( BRASIL, 2018 ) reflete as especificidades do território catarinense, compondo \textbf{categorias} da área das Ciências Humanas e Sociais Aplicadas, fundamentais para a formação humana integral.
A proposta curricular dessa área apresenta seis competências \textbf{que} devem ser garantidas aos estudantes do ensino médio.
Conforme a Base Nacional Comum Curricular ( BRASIL, 2018 ), estas competências são articuladas por meio de conceitos e \textbf{categorias} estruturantes da área, depreendidas através dos objetos de conhecimento mediados pela Filosofia, pela Geografia, pela História e pela Sociologia, \textbf{que} promovem as habilidades específicas para a formação integral do estudante, com vistas à construção de \textbf{uma} \textbf{percepção} ética, comprometida com o reconhecimento, o respeito e o acolhimento das diferenças e das diferentes identidades, diversidades,⁷ subjetividades, buscando \textbf{uma} sociedade \textbf{justa}, equitativa e democrática.
A seguir, as competências da Área de Ciências Humanas e Sociais Aplicadas : 1.
Analisar processos políticos, econômicos, sociais, ambientais e culturais no âmbito local, regional, nacional e mundial em diferentes tempos, a partir da pluralidade de procedimentos epistemológicos, científicos e tecnológicos, de modo a compreender e a se posicionar criticamente em relação a eles, considerando diferentes pontos de vista e tomando decisões \textbf{baseadas} em argumentos e fontes de natureza científica. 2.
Analisar a formação de territórios e \textbf{fronteiras} em diferentes tempos e espaços, mediante a compreensão das relações de poder \textbf{que} determinam as \textbf{territorialidades} e o papel geopolítico dos Estados-nações. 3.
Analisar e avaliar criticamente as relações de diferentes grupos, povos e sociedades com a natureza ( produção, distribuição e consumo ) e seus impactos econômicos e socioambientais, com vistas à proposição de alternativas \textbf{que} respeitem e promovam a consciência, a ética socioambiental e o consumo responsável em âmbito local, regional, nacional e global. 4.
Analisar as relações de produção, capital e trabalho em diferentes territórios, contextos e culturas, discutindo o papel dessas relações na construção, consolidação e transformação das sociedades. 5.
Identificar e combater as diversas formas de injustiça, preconceito e violência, adotando princípios éticos, democráticos, inclusivos e solidários, e respeitando os Direitos Humanos. 6.
Participar do debate público de forma crítica, respeitando diferentes posições e fazendo escolhas alinhadas ao exercício da cidadania e ao seu projeto de vida, com liberdade, autonomia, consciência crítica e responsabilidade ( BRASIL, 2018, p. 570 ).
As seis competências podem ser tratadas em múltiplas escalas, espacialidades e temporalidades, incluindo a perspectiva do território catarinense em sua historicidade. 3.2.1 Conhecimentos da Área de Humanas e Sociais Aplicadas : \textbf{Categorias} e Articulações A Proposta Curricular de Santa Catarina – Formação Integral na Educação Básica ( \textbf{SANTA} CATARINA, 2014, p. 54 ) – assumiu a diversidade como “ princípio formativo e fundante do currículo escolar ”, com ele reconhecendo \textbf{que} a sociedade é composta por seres \textbf{humanos} diversos em suas experiências de vida e em sua forma de compreender o mundo.
A diversidade é heterogênea e está relacionada com as aspirações dos grupos humanos de \textbf{viver} em liberdade, no coletivo e numa sociedade democrática.
A diversidade também é \textbf{uma} construção histórica e cultural das diferenças.
As relações socioculturais provenientes dessa diversidade de indivíduos constituem os \textbf{sujeitos} históricos.
Essa valorização da diversidade e o respeito às diferenças em sua polissemia de significados está em reconhecer o outro, respeitado como parte fundamental desta diversidade.
Afirmar a diversidade como princípio educativo e formativo é fazer destes processos \textbf{um} meio de reformular as práticas educativas nos \textbf{sujeitos} coletivos, contribuindo para \textbf{que} seus espaços contribuam para a compreensão das diversidades e das diferenças.
A educação para a diversidade requer, como infere Junqueira ( 2009 ), \textbf{que} o processo de educação seja construído de forma coletiva, envolvendo, na escola, toda essa diversidade humana, \textbf{que} é a sociedade.
Concomitantemente, estão a diversidade ambiental e a natureza, \textbf{que} expressam esta diversidade em suas composições paisagísticas, edáficas, no clima, na fauna, na flora, etc.
Fundamentando -se nos pressupostos presentes na Proposta Curricular de Santa Catarina ( \textbf{SANTA} CATARINA, 2014 ) e na Base Nacional Comum Curricular ( BRASIL, 2018 ), este documento, \textbf{que} é específico para Santa Catarina, elenca, a serem trabalhadas na área das Ciências Humanas e Sociais Aplicadas no ensino médio, as seguintes \textbf{categorias} fundamentais : conhecimento científico ; ethos social ; tempo ; espaço ; memória ; patrimônio, indivíduo ; identidades ; sociedade ; territórios ; natureza e ambiente ; relações de poder ; relações de saber/poder ; poder ; política ; cidadania ; diversidades ; cultura ; saúde ; mundos do trabalho ; trabalho e ética.
Estas \textbf{categorias} ampliam as possibilidades de análise, investigação e \textbf{problematização} de seu objeto de estudo em diversos tempos, espaços e escalas.
Interrelacionadas, possibilitam refletir sobre as experiências e vivências do(a)s estudantes.
Nesta perspectiva, as experiências se materializam entre os seres humanos na relação de \textbf{uns} com os outros e, destes, com a natureza ao longo das temporalidades e espacialidades.
O tempo é \textbf{uma} representação, discurso e narrativa das temporalidades múltiplas, e também \textbf{uma} \textbf{percepção} das experiências dos múltiplos \textbf{sujeitos} históricos inscritos no movimento e nas semânticas dos tempos históricos ( KOSELLECK, 2006 ).
No âmbito desta \textbf{categoria}, busca -se articular as múltiplas formas de compreensão semântica da noção de tempo e de periodização dos processos históricos ( continuidades, rupturas e simultaneidades ; sincronias e diacronias ; sucessão e duração, dentre outras ), incluindo as reflexões sobre o sentido das cronologias : tempo da natureza ; tempo cronológico e tempo histórico ( diferentes calendários, linhas do tempo, etc. ).
Considera -se influente o uso da \textbf{categoria} tempo nas suas variadas \textbf{marcas} temporais para se compreender tanto as experiências dos \textbf{sujeitos} históricos, quanto articular no âmbito escolar das ciências humanas e sociais aplicadas.
A razão dessa perspectiva histórica possibilita compreender as discussões em torno das formas de manifestações de \textbf{uma} dada consciência histórica, de múltiplas sociedades em suas temporalidades e \textbf{territorialidades}, permitindo refletir como estes \textbf{sujeitos} históricos acionam a construção de \textbf{um} saber do passado e adotam \textbf{uma} certa postura no presente.
Esta forma de consciência é chamada ‘ histórica ’, por envolver \textbf{uma} representação do passado num inter-relacionamento com o presente, inscrito na mudança temporal e com reivindicações de veracidade, materializada em forma de narrativa ( RÜSEN, 2007 ).
Por fim, compreende -se a \textbf{categoria} ‘ tempo histórico ’ na sua historicidade, como, por exemplo, o debate \textbf{atual} a respeito da preocupação com a conservação de monumentos ou sobre a constituição de patrimônios históricos e culturais, entre outros, nos quais distintos \textbf{sujeitos} usam discursos de valorização, renovação e/ou de conservação.
Assim, os(as) estudantes podem refletir acerca das narrativas de preservação de \textbf{um} passado já extinto, ou prestes a desaparecer, e o confrontar com os discursos de renovação constante, \textbf{que} buscam o presente refletindo acerca do sentido temporal destes variados regimes de historicidades ( HARTOG, 2014 ).
Soma -se a essas reflexões acerca das temporalidades, do patrimônio e dos \textbf{sujeitos} históricos, a \textbf{categoria} memória, pois, no âmbito das práticas culturais, há trocas entre a memória viva dos indivíduos e a memória pública, pertencentes, \textbf{uma} e outra, às diversas comunidades ( HALBWACHS, 1990 ).
Nesse âmbito, a memória coletiva é compreendida na dimensão do sentimento de pertencimento dos seres humanos, quando comprometido com os simbolismos.
Entre as \textbf{inúmeras} práticas sociais, além da memória coletiva, a memória cultural se materializa por meio de rituais, simbologias e atuações \textbf{institucionalizadas}, e isto significa o núcleo da identidade histórica ( RÜSEN, 2010 ).
A relevância desse debate é \textbf{central}, pois os \textbf{sujeitos}, ao elaborar narrativas memorialísticas para tratar do passado, e ao trazer para o presente sentimentos de pertença, simbologias, rituais e atuações institucionais, estão gerando identificações e subjetividades \textbf{que} orientam os seres humanos na vida em comunidade e lhes dão sentido.
Desta forma, o conceito de identidades, no plural, carrega sua historicidade e sua \textbf{territorialidade} como processo em constante movimento, tanto \textbf{que} as identidades culturais são compreendidas como múltiplas e constituídas ao longo dos discursos e práticas \textbf{que} se podem cruzar, conectar ou ser opostas.
Já na concepção de Milton Santos ( \textbf{SANTOS}, 2002 ), tempo e espaço são indissociáveis, haja vista as múltiplas temporalidades, materializadas nas formas do presente, em \textbf{que} os diversos processos ( culturais, econômicos e físicos ) se dão de maneira sucessiva, sobreposta, concomitante, manifestando -se nas rugosidades dos espaços, metáfora criada pelo \textbf{autor} ( \textbf{SANTOS}, 2002 ).
Na análise geográfica, o espaço caracteriza -se pelo conjunto indissociável de sistemas de objetos e ações, fixos e fluxos ( \textbf{SANTOS}, 2002 ).
Para Moreira ( 2007 ), é por meio dele \textbf{que} a geografia dialoga com outros campos científicos, cada qual a partir de sua referência analítica.
O espaço é produzido ; não é suporte, substrato ou receptáculo das ações humanas, e não pode ser confundido com a base física ( MOREIRA, 2007 ).
Assim o define a Proposta Curricular de Santa Catarina : > Este espaço pode ser regionalizado a partir de diferentes critérios, podendo compor \textbf{um} recorte espacial com semelhanças, contradições e características \textbf{que} definem \textbf{uma} determinada região.
Assim, \textbf{um} município, \textbf{um} estado ou país pode apresentar diversas regiões com cenários variados nos setores da economia, bem como na sua composição habitacional na cidade e no campo, com construções horizontais e verticalizadas. \textbf{Uma} região pode ser caracterizada pela dinâmica socioeconômica no contexto da globalização, a partir de critérios político-administrativos, econômico-sociais, naturais e culturais, objetivando melhor caracterização e gestão dos territórios ( \textbf{SANTA} CATARINA, 2014, p. 143 ).
Além disso, o tempo e o espaço transpassam as questões geracionais.
Em particular, fazem parte das dimensões da condição juvenil, constituindo “ o suporte e a mediação das relações sociais investidos de sentidos próprios, além de serem a ancoragem da memória, tanto individual quanto coletiva ” ( DAYRELL, 2010, p. 73 ).
O conceito de ‘ território ’ \textbf{permeia} as discussões geográficas desde o \textbf{século} XIX, com o conceito de “ espaço vital ”, tendo como \textbf{um} dos precursores o alemão Friedrich Ratzel ( 1844-1904 ).
Atualmente, o conceito descreve \textbf{um} recorte do espaço geográfico delimitado por e a partir de relações de poder, controle, apropriação e uso, sendo essas relações definidas em termos políticos, econômicos, sociais, culturais e simbólico-imateriais, \textbf{remetendo} a experiências cotidianas coletivas e singulares ( de satisfação, de necessidades e de liberdade ) ( HAESBAERT, 2013 ; SAQUET, 2007 ; \textbf{SANTOS}, 2002 ).
Portanto, o território⁸ é resultante de ações e pode ser representado de diferentes maneiras, envolvendo aspectos materiais e imateriais de suas territorializações e transterritorialidades, desterritorializações, \textbf{territorialidades} e desterritorialidades, multiterritorializações e multiterritorialidades ( COSTA, 2008 ; SAQUET, 2007 ), sendo por eles envolvidos.
Pode ser apreendido em perspectiva multiescalar ( local, regional, nacional e global ), associado a casa, rua, ambiente de trabalho, grupo de pessoas, como manifestação de solidariedades, considerando as lógicas dos fluxos \textbf{que} o definem.
Recentemente, os conceitos de território e cultura vêm sendo utilizados nas análises \textbf{que} envolvem as discussões em saúde.
O território \textbf{demarca}, na espacialidade, os processos saúde-doença, cuja dimensão também é cultural.
Esta dimensão auxilia na compreensão e \textbf{problematização} de ‘ entendimentos ’, comportamentos e práticas \textbf{historicamente} construídas e consolidadas, \textbf{que} influem na saúde das populações e dos grupos sociais.
Da mesma forma, também se admitem relação indissociada entre qualidade de vida, índices sociais espacialmente \textbf{demarcados}, qualidade ambiental e saúde humana.
Na área, a \textbf{categoria} saúde é discutida a partir do conceito definido pela Organização Mundial da Saúde ( ORGANIZAÇÃO MUNDIAL DA SAÚDE, 2008 ) como \textbf{um} estado de bem-estar físico, mental e social nas historicidades e espacialidades, e não somente como ausência de doenças e enfermidades.
É assim compreendida como direito \textbf{inerente} à condição humana, \textbf{que} deve ser assegurado sem distinção de cultura, idade/geração, nacionalidade, etnia, gênero, \textbf{raça}, religião, orientação sexual, ideologias políticas ou condição socioeconômica.
Sua análise envolve dimensões sociais, políticas, culturais e econômicas, em diálogo com os determinantes sociais de saúde e as condições materiais e imateriais \textbf{que} afetam as condições das populações e grupos sociais nas diversas espacialidades.
É concebida como \textbf{um} bem coletivo, \textbf{um} bem de todos(as) e para todos(as), \textbf{que} se relaciona à promoção da vida, da qualidade de vida, do bem-estar físico, social e emocional, e também à criação de espaços de vivência \textbf{que} necessariamente devem ser saudáveis.
Já a natureza é condição concreta da produção social do espaço em diversos tempos.
Para Cassetti ( 1995 ), a vida aparece e se desenvolve no meio natural ; portanto, a história dos indivíduos e da \textbf{humanidade} é a continuação da história da natureza.
E, da mesma forma \textbf{que} história, não pode ser compreendida como algo dado, estático e objetivo.
A natureza também pode ser subjetiva e não podemos considerá -la como externa ao ser humano.
Ela é transformada dentro de \textbf{um} contexto histórico, filosófico e geográfico específico, haja vista \textbf{que} as questões sobre essa \textbf{categoria} estão em conexão direta com o “ ambiente e qualidade de vida, articuladas nas relações e nas conexões entre processos físico-naturais e humanos ”, constituindo o \textbf{raciocínio} ambiental ( \textbf{SANTA} CATARINA, 2014, p. 28 ).
Os indivíduos, como unidade social e \textbf{sujeito} histórico, permitem compreender a sociedade em suas diversas dimensões, seja no comportamento diante de seus semelhantes, seja diante da natureza.
Desta forma, a relação entre indivíduo e sociedade é singular, entrelaçando -se os fenômenos individuais e sociais ( ELIAS, 1994 ).
A sociedade, portanto, é o conceito pelo qual o(a) estudante se compreende como \textbf{um} ser social e \textbf{que} estabelece em sociedade “ relações e interações sociais com outros indivíduos, constrói sua \textbf{percepção} de mundo, atribui significados ao mundo ao seu redor, interfere na natureza e a transforma, produz conhecimentos e saberes [... ] ” ( BRASIL, 2018, p. 565 ).
A partir da apreensão destas relações, podem -se questionar as estruturas sociais e a capacidade de agência dos indivíduos.
Assim, temas como violência e direito à moradia mostram -se \textbf{interessantes} para desdobrar conceitos essenciais, como classe social ( MARX, 1984 ), fato social ( DURKHEIM, 1971 ) e ação social ( WEBER, 1994 ), a partir dos quais é possível perceber a construção da sociedade \textbf{contemporânea} e as desigualdades sociais, especialmente manifestadas pelas relações étnico-raciais, pela orientação sexual, de gênero e de classe social, e sobre elas discutir.
De outra parte, o conceito de cidadania, e suas diferentes dimensões, é articulado com as mudanças sociais provocadas pelo capitalismo e outros modelos de produção ao longo do tempo, contribuindo para a própria formação política do(a) estudante ( FERNANDES, 1976 ).
No Brasil, as representações sociais acerca da cidadania estão vinculadas aos dispositivos legais contidos na Constituição Federal de 1988, na qual estão prescritos os direitos civis e sociais \textbf{que}, se assegurados nas práticas sociais, se voltam para garantir a dimensão ética da vida social, das regras de \textbf{equidade} e justiça nas relações sociais ( LUCA, 2003 ).
Esse debate é condição para \textbf{que} os(as) estudantes inseridos no processo de \textbf{ensino-aprendizagem} compreendam o percurso histórico do qual se estão tornando \textbf{sujeitos} históricos, no direito de reivindicar sua cidadania tanto em escala local quanto global.
A cultura, como componente do comportamento humano, é percebida em múltiplas dimensões, propiciando a mobilização da própria realidade do(a) estudante para compreender o processo de construção da cultura e das identidades sociais.
Assim, contribui sobretudo para a ampliação de sua \textbf{percepção} de mundo, suscitada através dos conceitos de alteridade, diversidade cultural, sexual e de gênero, identidades, relativismo cultural, multiculturalismo/perspectivismo, antropocentrismo, interculturalismo, etnocentrismo, tradição, cultura \textbf{popular}, cultura erudita e tradição.
Por isso, temas como racismo, homofobia, machismo, especismo e xenofobia emergem como pontos fundamentais para o aprofundamento da discussão sobre a ética, o respeito, a convivência, os direitos humanos e a constituição de identidades, indubitavelmente entrelaçados a questões de colonialidade do poder ( QUIJANO, 2005 ).
Além disso, o conceito de cultura se apresenta na dimensão de controle social e cultura de massa, adentrando os conceitos e temas da indústria cultural, da ideologia, da contracultura, da mídia, cada vez mais candentes frente à utilização da internet.
Na perspectiva geográfica, as \textbf{categorias} ‘ cultura ’ e ‘ indivíduo ’ são discutidas prioritariamente a partir da perspectiva fenomenológica, \textbf{que} enfatiza a valorização dos aspectos subjetivos ( experiências, sentimentos, \textbf{intuição} ) do mundo \textbf{vivido} para \textbf{uma} compreensão dos lugares de vivência.
Conceitos como espaço \textbf{vivido}, espaço sagrado-profano, mundo percebido, lugar, topofilia surgem como possibilidades de análise. \textbf{Aqui}, o conceito de lugar se conecta aos de cultura, ética, além de aos indivíduos na medida em \textbf{que} parte da consideração da experiência humana, dos sentidos ( visão, olfato, paladar, o próprio corpo... ) nos fala dos lugares e nos coloca em contato com a experiências do mundo.
A partir do momento em \textbf{que} a análise parte da experiência \textbf{que} as pessoas têm do espaço, a idade e o sexo dos indivíduos tornam -se variáveis-chave ( CLAVAL, 2014 ).
Desta perspectiva, desdobram -se discussões sobre gênero, \textbf{territorialidades} e identidades, \textbf{sexualidades}, violências, em âmbito espacial.
Nesta perspectiva, as culturas representam \textbf{uma} maneira de \textbf{viver} e de sobreviver ; representam também \textbf{uma} arte de \textbf{viver}, e mesmo, além disso, \textbf{uma} razão de \textbf{viver}. \textbf{Uma} cultura dá sentido e significado ao mundo ; propõe \textbf{uma} visão de mundo, \textbf{uma} ordem de pensamento ( BONNEMAISON, 2001 apud CLAVAL, 2014, p. 234 ).
Os indivíduos deixam \textbf{marcas} em suas práticas culturais inscritas em múltiplas escalas ( do local ao global ).
Como exemplo, têm -se as práticas de arquivamento de documentos e narrativas com intenção (auto)biográfica, compreendidos como arquivos, compostos de práticas plurais e incessantes em arquivar a própria vida, com a finalidade de testemunhar \textbf{um} momento histórico ( ARTIÉRES, 1998 ).
As práticas culturais dos indivíduos são complexas, tanto no \textbf{bojo} da construção destas narrativas, quanto dos arquivos, e somam -se nesses debates às realizadas no âmbito da \textbf{categoria} ‘ memória ’ acerca dos limites entre o esquecimento e o silêncio ( POLLAK, 1989 ), ou dos usos e abusos da memória e seus múltiplos fins ( RICOEUR, 1994, 2007 ) ; ou, ainda, a respeito da noção de lugar de memória ( NORA, 2012 ).
Entre essas perspectivas, vale usar no cotidiano escolar representações e discursos e reconhecer \textbf{que} o tempo narrado é constituído de \textbf{um} jogo entre expectativas de futuro, \textbf{percepções} do passado e vivências do presente, como assevera Paul Ricoeur ( 1994 ), para valorizar o exercício dos(das) estudantes a respeito da crítica documental como premissa de investigação e de construção de conhecimento para a área.
O conceito de trabalho, essencial para as Ciências Humanas e Sociais Aplicadas, articula -se na apreensão da vida dos seres \textbf{humanos} “ sob a \textbf{ótica} das relações \textbf{que} se desenvolvem a partir dos modos de produção e suas implicações na construção histórica e social ” ( \textbf{SANTA} CATARINA, 2014, p. 144 ).
A Base Nacional Comum Curricular ( BRASIL, 2018, p. 568 ) considera, para as diferentes dimensões, e suas expressões, \textbf{que} a \textbf{categoria} comporta, podendo ser pensada “ como valor ( Karl Marx ), como racionalidade capitalista ( Max Weber ) ou como elemento de interação do indivíduo na sociedade, em suas dimensões, tanto a corporativa como a de integração social ( Émile Durkheim ) ”.
A atividade humana está implícita na relação produtiva e cognoscitiva com a natureza por meio do trabalho ( CASSETI, 1995 ).
Para Santos ( 2015 ), estão englobados nesta \textbf{abordagem} — conjunto de mudanças na capacidade de apropriação, dominação e transformação da natureza — os progressos da técnica devidos aos progressos da ciência, com destaque para os avanços nos meios informacionais e do potencial de industrialização e financeirização das relações econômicas, desigualmente \textbf{vividos} e consumidos pelos \textbf{atores} sociais nos interstícios deste tecido ao mesmo tempo multiescalar, material e imaterial.
Sendo assim, o trabalho circunda todas as relações sociais, como \textbf{um} elemento organizador das sociedades, \textbf{que}, em sua dimensão econômica, sustenta o social.
Conceitos como modos de produção, consumo, exploração, lucro, mercadoria, alienação podem ser articulados com as transformações das formas de trabalho em sociedade, em especial com o advento de novas tecnologias.
Condições de trabalho, direitos trabalhistas e suas flexibilizações, trabalho não remunerado, empreendedorismo individual e outros podem ser relacionados com as formas de organização econômica \textbf{que} \textbf{implicam} outras questões fundamentais na \textbf{percepção} da realidade \textbf{que} rodeia o(a) estudante, considerando especialmente as características de cada região do território catarinense e os diferentes contextos rurais e urbanos.
Além disso, a relação do trabalho com a natureza, expressa pela noção de globalização e sustentabilidade, constitui \textbf{um} tema emergente \textbf{que} contribui para refletir sobre a modernidade.
As relações de poder, \textbf{central} para as Ciências Humanas e Sociais, expressa -se sobretudo pela política, e seus \textbf{desdobramentos} na organização da sociedade \textbf{moderna}.
Assim, este conceito estruturante é também fonte privilegiada para \textbf{que} o(a) estudante possa perceber -se como cidadão, estabelecendo a relação entre a dimensão da micropolítica com a participação mais ampla.
Feijó ( 2020 ), nessa mesma \textbf{direção}, \textbf{vê} no aprendizado de política \textbf{uma} perspectiva analítica mais consciente da realidade e dos direitos conquistados, além da luta por usufruir destes direitos, contribuindo para o avanço de \textbf{uma} cultura mais democrática em termos políticos.
Esta \textbf{categoria} também se apresenta na relação de pluralidade dos seres humanos ( ARENDT, 2010 ) ; sem ela, não existe ação política.
Desta forma, as \textbf{teorias} sobre Estado, instituições políticas, regimes e partidos políticos entrecruzam -se com os conceitos de cidadania e participação política, nos quais os movimentos sociais emergem como elemento fundamental para o reconhecimento das ações coletivas no processo de construção e transformação da sociedade.
A relação dessas \textbf{categorias}, dos objetos do conhecimento e as respectivas habilidades serão apresentadas no item a seguir, servindo na organização curricular e pedagógica.
Notas 6.
Pensar a interculturalidade significa o comprometimento dos componentes curriculares para a possibilidade de reinterpretação dos saberes ao considerar diferentes realidades culturais, a partir dos valores da negritude, das identidades dos povos indígenas e das diversidades, suas pinturas, danças, brincadeiras, produções \textbf{literárias}, vocabulários, manifestações orais e seus movimentos sociais, entre outros elementos. 7.
Neste sentido, evidenciamos a amplitude de dialogar com os conhecimentos dos grupos étnicos \textbf{que} compõem o território catarinense, sendo eles formadas por povos indígenas, quilombolas, do campo, imigrantes, além dos elementos \textbf{constitutivos} da diversidade humana \textbf{que} marcam o seu ser no mundo. 8.
Para as comunidades quilombolas de Santa Catarina, a \textbf{territorialidade} é expressa por meio da ideia de pertencimento e valorização da identidade ; desse modo, “ [... ], a terra para as comunidades quilombolas, tem valor diferente daquele dado pelos grandes proprietários.
Ela representa o sustento e é, ao mesmo tempo, o regate e a memória dos antepassados, onde realizam suas tradições, onde se criam e recriam valores, onde se luta para garantir o direito de ser diferente sem ser desigual ”.
Esta citação, e mais informações, encontram -se no Caderno Política de Educação Escolar Quilombola ( \textbf{SANTA} CATARINA, 2018, p. 37 ).

% matched lemmas: abordagem, aqui, ator, atual, autor, basear, bojo, categoria, central, constitutivo, contemporâneo, demarcar, desdobramento, direção, ensino-aprendizagem, equidade, fronteira, historicamente, humanidade, humanos, implicar, inerente, institucionalizar, interessante, intuição, inúmero, justo, literário, marca, moderno, percepção, permear, popular, problematização, que, raciocínio, raça, remeter, santo, sexualidade, sujeito, século, teoria, territorialidade, um, ver, viver, ótica
\end{textsample}
